\chapterwithopening{Strategic Vision and Project Scope}{chap:strategic_vision}{
  \chaptercontentsitem{sec:project_vision}
  \chaptercontentsitem{sec:problem_statement}
  \chaptercontentsitem{sec:project_motivations}
  \chaptercontentsitem{sec:strategic_objectives}
  \chaptercontentsitem{sec:team_goals}
  \chaptercontentsitem{sec:project_scope}
}

\section{Project Vision}\label{sec:project_vision}

Digital transformation has become a defining condition of contemporary
education. Institutions increasingly depend on online systems not only to store
materials, but also to coordinate communication, teaching, assessment,
feedback, collaboration, and daily academic follow-up across time and location.
As that dependence grows, the quality of the learning environment depends less
on the presence of isolated tools and more on whether those tools form a
coherent academic system.

Within this broader context, educational workflows often remain fragmented.
Class communication may occur in one platform, materials in another, live
teaching in another, and practical technical work elsewhere. EduVerse is
proposed as a response to that condition: an integrated educational platform
that brings related academic workflows into one governed digital environment
instead of leaving them scattered across loosely connected services.

This strategic direction is intentionally broader than any single institutional
type. EduVerse is conceived as a governed academic environment that can serve
universities, schools, online academies, bootcamps, tutoring centers,
exam-preparation spaces, corporate training contexts, and open-learning
communities without assuming that all of them need the same configuration.

Its vision is shaped by two practical goals. The first is to provide a
coherent digital environment in which related academic work can remain
connected. The second is to preserve adaptability, so that organizations and
classes can enable, disable, extend, or restrict features according to their
own policies and teaching needs. EduVerse is therefore framed not as a fixed
bundle of tools, but as a configurable academic workspace designed for
continuity, governance, and long-term growth.

\section{Problem Statement}\label{sec:problem_statement}

Many educational workflows are currently distributed across separate systems.
One platform may be used for class communication and sharing materials, another
for live meetings and presentations, another for collaborative whiteboard
interaction, and another for programming practice. Familiar examples include
Google Classroom for announcements and materials, Google Meet for live calls,
Excalidraw-like tools for shared whiteboard work, and LeetCode-like
environments for coding exercises. While each tool can be useful on its own,
their separation introduces repeated context switching, inconsistent workflow
boundaries, weaker continuity between related academic tasks, and scattered
responsibility for governance and follow-up.

This fragmentation becomes more serious in multi-organization educational
settings. A platform must not only expose features; it must also manage roles,
permissions, memberships, public access rules, class visibility, and the
relationship between administrative governance and daily learning activity.
When these concerns are spread across unrelated services, both usability and
accountability become harder to maintain.

The problem also has a client-design dimension. Rich administrative workflows,
content management, assessment authoring, and complex class control are better
suited to a full web environment. At the same time, many high-frequency actions
such as checking updates, opening materials, reading chat, reviewing class
summaries, and joining live sessions should remain practical from mobile
devices. A realistic educational platform therefore needs both a strong primary
web application and a carefully scoped mobile companion rather than an
unstructured duplication of screens.

\section{Project Motivations}\label{sec:project_motivations}

EduVerse was motivated by the opportunity to build a digital learning
environment in which academic continuity and institutional control could be
designed together. Rather than treating governance, communication, assessment,
and learning support as loosely related concerns, the project aimed to realize
them within one environment that could remain useful across different
educational contexts.

The project was also academically and professionally valuable for the team. It
required the combination of requirements analysis, system modeling,
implementation, service integration, security-aware design, testing, and
technical reporting within one sustained engineering effort. In this sense,
EduVerse served not only as a software product, but also as a realistic
graduation project through which the team could practice disciplined
collaborative development.

\section{Strategic Objectives}\label{sec:strategic_objectives}

The strategic objectives of EduVerse translate the project vision into concrete
institutional and functional targets.

\begin{itemize}
  \item To provide a multi-organization educational platform that can support
    distinct institutions or organizational units within one governed system.
  \item To support role-based operation for Organization Administrators,
    Teachers, and Students, including users who belong to more than one
    organization or hold more than one role.
  \item To unify classes, materials, chat, assignments, exams, results,
    notifications, live sessions, whiteboard collaboration, AI support, and
    coding practice within one academic environment.
  \item To allow supported features to be enabled or disabled at organization
    level and, where appropriate, refined further at class level.
  \item To allow AI-supported features to be disabled for organizations or
    classes that do not wish to use them.
  \item To support public or open-learning access where enabled, so that
    organizations can share educational materials or membership paths more
    broadly without giving up institutional control.
  \item To support controlled extension and future integration so that new
    capabilities or external clients can reuse the same protected academic API
    model rather than rebuilding business rules from scratch.
  \item To provide a primary web platform for complete academic and
    administrative workflows while extending selected day-to-day actions
    through a mobile companion application.
\end{itemize}

\section{Educational and Professional Goals for the Team}\label{sec:team_goals}

For the project team, EduVerse represents a complete software engineering
exercise grounded in a realistic problem domain. The work required the team to
move from problem framing to requirements analysis, methodology selection,
system design, implementation, validation, and documentation in a coordinated
way.

The project also supported professional growth. It strengthened experience in
problem decomposition, collaborative planning, technical trade-off analysis,
review discipline, and responsibility for system quality. These outcomes are
important because the value of a graduation project lies not only in the final
software, but also in the maturity of the engineering process used to produce
it.

\section{Project Scope and Boundaries}\label{sec:project_scope}
\subsection{System Boundaries and User Roles}\label{subsec:system_boundaries_roles}

EduVerse operates through three principal human roles: Organization
Administrator, Teacher, and Student. The same person may hold more than one
role, and a user may belong to more than one organization. For that reason, the
system resolves access according to the selected organization, the active role,
class membership, and the current feature configuration rather than by assuming
that every user belongs to one fixed category.

This role structure defines an important project boundary. EduVerse is not a
general social platform in which all users share the same permissions. It is a
governed educational workspace in which visibility, management authority,
assessment control, and feature access are deliberately differentiated.

\subsection{Functional Scope}\label{subsec:functional_scope}

Within its implemented scope, EduVerse includes authentication, organization
access and membership, class management, materials, chat, assignments, exams,
results, notifications, live sessions, whiteboard collaboration, an AI Agent,
and an integrated coding workspace. It also includes organization-level and
class-level feature controls, extension support, and optional public or
open-learning access where those capabilities are enabled.

This scope is not restricted to higher education alone. The same model can be
adapted for primary schools and kindergartens through narrower feature
selection and closer organization-level control, just as it can support more
specialized contexts such as technical courses, tutoring programs, or open
learning initiatives.

The class model is intentionally flexible. A class can support a conventional
course, but it can also be used for a narrower academic purpose. For example, a
class may function as a digital library space dedicated mainly to sharing books
and learning materials. This flexibility shows that EduVerse is not restricted
to one teaching pattern; it is a configurable academic workspace that can be
adapted to different institutional uses.

\subsection{Web and Mobile Client Scope}\label{subsec:web_mobile_client_scope}

The EduVerse web application is the primary and most complete client. It
supports the broader administrative, instructional, and assessment workflows
that define the platform as a full educational system.

EduVerse Mobile already exists as a companion mobile application, developed to
extend selected platform features to mobile users. Its scope is intentionally
narrower than the web platform. The current mobile implementation supports
authenticated access, class summaries, notifications, materials, class chat,
organization switching, and basic live-session participation. It is therefore
best described as a companion under active development rather than as a
feature-complete replacement for the web platform. Administrative workflows and
several advanced learning features remain mainly web-side.

\subsection{Distinctive Features and Extensible Design}
\label{subsec:competitive_features_innovation}

EduVerse is distinctive because it combines integration with control. The
platform brings together communication, materials, live collaboration,
assessment, AI assistance, and coding practice, but it does so inside a
governed multi-organization model rather than through a loose collection of
independent tools.

Its extensible design is also central to the project. Supported features can be
enabled or disabled at organization level and, where supported, refined further
at class level. AI features can likewise be disabled when an organization or a
specific class considers them unsuitable. The same underlying API model also
creates reuse potential for future integrations or companion clients, because
core workflows are exposed through protected shared routes rather than being
hard-wired to one interface only.

Public access adds another important dimension. When enabled, public
organizations or public join paths can support broader sharing of educational
materials and more open learning participation, while still remaining under
organization-defined governance rules.
